\section*{Resumen}

Paralegales es una plataforma de \textit{legaltech} orientada a automatizar el seguimiento de procesos judiciales en la Rama Judicial Unificada en Colombia. El objetivo principal de esta solución es simplificar la tarea de los abogados al ofrecer herramientas que permiten monitorear automáticamente los cambios en los procesos judiciales, además de proporcionar funcionalidades complementarias para mejorar la eficiencia en la gestión legal.

Para acelerar la salida a producción, el desarrollo inicial de Paralegales se apoyó en el \textit{Software Development Kit} (SDK) de Firebase, lo que permitió avanzar rápidamente en funcionalidades clave. Sin embargo, esta decisión acarreó deuda técnica significativa: lógica de negocio crítica implementada directamente en el \textit{frontend} (con los riesgos de seguridad asociados al patrón \textit{Insecure Design} del OWASP Top 10), fragmentación arquitectónica entre funciones AWS Lambda y servicios Firebase sin patrones claros que gobernaran sus fronteras, ausencia de observabilidad transversal sobre las integraciones del sistema y carencia de prácticas formales de aseguramiento de calidad en el ciclo de desarrollo.

El presente proyecto, en lugar de acometer una migración total hacia AWS —costosa, disruptiva para un producto ya en producción y descartada por la organización tras evaluar el caso de la migración del motor de persistencia—, adopta una estrategia de \textbf{evolución arquitectónica híbrida}. Esta estrategia se materializa en cuatro frentes complementarios: (i) la introducción de un patrón \textit{Backend-for-Frontend} (BFF) que reubica la lógica sensible fuera del cliente; (ii) la refactorización de la arquitectura interna del \textit{frontend} bajo principios de \textit{Domain-Driven Design} (DDD) con Nuxt Layers, y del \textit{backend} bajo una separación estricta \textit{Handler--Service--Repository} que preserva la opción de evolución futura del motor de persistencia; (iii) la construcción de un sistema transversal de observabilidad estructurada apoyado en interceptores HTTP, CloudWatch y notificaciones proactivas; y (iv) la consolidación de prácticas continuas de aseguramiento de calidad mediante \textit{hooks} de Git, integración continua, pruebas \textit{end-to-end} y pruebas unitarias de reglas de seguridad.

El resultado es una plataforma con menor superficie de exposición de lógica crítica, mejor trazabilidad operativa, mayor mantenibilidad y preparada para una evolución incremental controlada, sin comprometer la continuidad del servicio ni acoplarse irreversiblemente a un único proveedor de nube.

\hfill

\textbf{Palabras clave:} Ingeniería de software, Arquitectura híbrida, \textit{Backend-for-Frontend}, \textit{Domain-Driven Design}, Microservicios \textit{serverless}, Observabilidad, Aseguramiento de calidad, \textit{Legaltech}.
